\documentclass[a4paper,11pt]{article}
\usepackage[italian]{babel}
\usepackage[margin=18mm]{geometry}
\usepackage{array,booktabs,tabularx,xcolor}
\usepackage[hidelinks]{hyperref}

\definecolor{fixdblue}{HTML}{1E3A52}
\newcolumntype{L}[1]{>{\raggedright\arraybackslash}p{#1}}
\newcolumntype{C}[1]{>{\centering\arraybackslash}p{#1}}
\newcolumntype{Y}{>{\raggedright\arraybackslash}X}
\setlength{\parindent}{0pt}
\setlength{\parskip}{5pt}

\begin{document}

\begin{center}
  {\LARGE\bfseries\color{fixdblue} Priorità di implementazione di Fixd\par}

\end{center}

La prima versione utilizzabile (MVP) permette di cercare un servizio,
consultarne le disponibilità e contattare il professionista dopo l'accesso.
La ricerca e la consultazione sono pubbliche; pubblicazione, modifica dei
propri annunci e contatto richiedono un account autenticato.

Le colonne seguono il modello del foglio Excel. L'ordine di sviluppo va
 da 1 a 8; l'importanza è espressa su una scala da 1 (minima) a 10
 (massima). Le stime sono in giorni-persona.
 L'ultima colonna indica l'attore che avvia la funzione.

\medskip
{\small
\setlength{\tabcolsep}{3pt}
\renewcommand{\arraystretch}{1.2}
\begin{tabularx}{\linewidth}{@{}YC{1.7cm}C{2.1cm}C{2.1cm}L{4.0cm}@{}}
  \toprule
  \textbf{Funzionalità generiche} & \textbf{Ordine di sviluppo} &
  \textbf{Importanza presenza funzione} & \textbf{Tempo previsto (giorni)} &
  \textbf{Funzione attivata da} \\
  \midrule
  \textbf{Registrazione} & 1 & 9 & 4 & Utente non autenticato \\
  \addlinespace[5pt]
  \textbf{Accesso e autorizzazioni} & 2 & 10 & 4 & Utente non autenticato \\
  \addlinespace[5pt]
  \textbf{Gestione annunci} & 3 & 10 & 7 & Utente autenticato (autore) \\
  \addlinespace[5pt]
  \textbf{Ricerca annunci} & 4 & 10 & 9 & Utente, anche non autenticato \\
  \addlinespace[5pt]
  \textbf{Filtro per categoria} & 5 & 8 & 3 & Utente, anche non autenticato \\
  \addlinespace[5pt]
  \textbf{Contatto del professionista} & 6 & 10 & 4 & Utente autenticato \\
  \addlinespace[5pt]
  \textbf{Calendario delle disponibilità} & 7 & 9 & 7 & Consultazione: ogni utente. Gestione: autore autenticato \\
  \addlinespace[5pt]
  \textbf{Ricerca assistita dall'AI} & 8 & 5 & 9 & Utente, anche non autenticato \\
  \bottomrule
\end{tabularx}
\par}

\medskip


\end{document}
